(1) 径向方程为 \(m\ddot r=\ell^2/(mr^3)-kr^{-p}\)，圆轨道 \(r_c\) 需满足 \(\ell^2/(mr_c^3)=kr_c^{-p}\)，即 \(kr_c^{3-p}=\ell^2/m\)。在 \(p=3\) 的退化情况下不能用这个关系单独确定稳定的半径。
(2) 径向有效势满足 \(U_{\rm eff}'(r)=-\ell^2/(mr^3)+kr^{-p}\)。二阶导在圆轨道处为 \(U_{\rm eff}''(r_c)=[3-p]kr_c^{-p-1}\)。由于 \(k,r_c>0\)，小扰动具有恢复力的充分必要条件为 \(p<3\)；\(p=3\) 是临界退化。
(3) 线性化运动 \(m\ddot{\delta r}+U_{\rm eff}''(r_c)\delta r=0\)，故 \(\omega_r^2=(3-p)kr_c^{-p-1}/m\)。又由向心力平衡 \(\Omega^2=\ell^2/(m^2r_c^4)=kr_c^{-p-1}/m\)，两式相除得到 \(\omega_r/\Omega=\sqrt{3-p}\)。
(4) 两个相邻近心点隔一个径向周期 \(2\pi/\omega_r\)，对应方位角 \(\Delta\phi=2\pi\Omega/\omega_r=2\pi/\sqrt{3-p}\)。\(p=2\) 给 \(\Delta\phi=2\pi\)，\(p=-1\) 给 \(\Delta\phi=\pi\)，分别对应开普勒近圆轨道和各向同性谐振子轨道的特征。
